\thispagestyle{empty}\vspace*{\fill}
\begin{minipage}{135mm}\fontsize{9.5}{14}\selectfont\color{otMuted}
{\otSemi\color{otInk}Physique-chimie · Seconde générale et technologique}\par\medskip
Collection OptiTAB. Édition du 25 octobre 2026.\par
240 exercices de chapitre, 20 défis et six évaluations originales corrigées.\par\medskip
Référence : programme français de physique-chimie de Seconde générale et technologique,
BO spécial n° 1 du 22 janvier 2019, arrêté du 17 janvier 2019, NOR MENE1901634A.
Programme signalé en vigueur sur Éduscol lors de la vérification du 25 octobre 2026.
L'ordre des chapitres et les formats d'évaluation sont des choix de la collection.\par\medskip
Ce manuel indépendant n'est pas une publication du ministère de l'Éducation nationale.
Les cours, situations d'entraînement et figures ont été rédigés pour cet ouvrage ; aucune
annale ni aucun exercice d'éditeur n'est reproduit. Les mesures des exercices sont construites
ou idéalisées pour l'entraînement. Les constantes et les références sont indiquées dans les annexes.\par\medskip
Les schémas sont vectoriels. Les échelles utiles à une construction ou une lecture sont précisées ;
les autres dessins sont des schémas de principe. Les activités expérimentales se réalisent
avec l'encadrement et le matériel adaptés au laboratoire scolaire.\par\medskip
© OptiTAB, 2026.\quad\otSite
\end{minipage}\clearpage
